
\subsection{Automatisiertes Zurücksetzen der LDAP-Testumgebung}

Da durch die Ausführung der Testfälle Benutzer im LDAP-System angelegt, aktualisiert oder gelöscht werden, musste nach jedem Testlauf ein definierter Ausgangszustand wiederhergestellt werden. Um diesen Prozess zu vereinfachen, wurde ein Bash-Skript entwickelt, welches die Testumgebung automatisiert neu aufsetzt.

\subsubsection*{Zielsetzung}

Das Skript stellt sicher, dass:

\begin{itemize}
  \item der OpenLDAP-Container neu gestartet wird,
  \item alle relevanten LDIF-Dateien neu importiert werden (Schema, Struktur, Gruppen, Benutzer),
  \item eine konsistente Ausgangsbasis für jeden Testlauf garantiert ist.
\end{itemize}


\subsection{Automatisiertes Zurücksetzen der LDAP-Testumgebung}

Da durch die Ausführung der Testfälle Benutzer im LDAP-System angelegt, aktualisiert oder gelöscht werden, musste nach jedem Testlauf ein definierter Ausgangszustand wiederhergestellt werden. Um diesen Prozess zu vereinfachen, wurde ein Bash-Skript verwendet, das die Testumgebung automatisiert neu aufsetzt.

\subsubsection*{Herkunft und Motivation}

Das Skript wurde mit Hilfe von \textbf{ChatGPT (OpenAI)} generiert und anschlieend in das Projekt übernommen. Es basiert funktional auf den Anweisungen im \texttt{README.md} der Entwicklungsumgebung, in dem ursprünglich alle Schritte manuell beschrieben waren - wie z.\,B. das Stoppen und Starten der Container sowie der Import von Schema- und Beispieldaten. Ziel war es, diese manuellen Kommandos in ein einziges, robustes Skript zu überführen.

\subsubsection*{Zielsetzung}

Das Skript stellt sicher, dass:

\begin{itemize}
  \item der OpenLDAP-Container vollständig neu gestartet wird,
  \item die im Projekt enthaltenen LDIF-Dateien korrekt geladen werden (\texttt{schema.ldif}, \texttt{01\_ou\_structure.ldif}, \texttt{02\_groups.ldif}, \texttt{03\_users.ldif}),
  \item eine konsistente und saubere Ausgangsbasis für jeden Testlauf gegeben ist.
\end{itemize}

\subsubsection*{Skriptübersicht}

\begin{lstlisting}[language=bash, caption={Von ChatGPT generiertes Skript - Automatisierung der README-Kommandos}]
#!/bin/bash

set -e
set -o pipefail

# Git Bash Workaround
export MSYS_NO_PATHCONV=1

echo ">>> Stopping and removing containers and volumes..."
docker compose down -v

echo ">>> Starting containers..."
docker compose up -d

echo ">>> Waiting for OpenLDAP container to be ready..."
sleep 5

# Check container
if ! docker ps --format '{{.Names}}' | grep -q '^openldap$'; then
  echo " Error: OpenLDAP container 'openldap' not running."
  exit 1
fi

# Validate LDIF files
docker exec openldap ls -la /container/service/slapd/assets

# Import LDIF files
winpty docker exec openldap sh -c "ldapadd -Y EXTERNAL -H ldapi:/// -f /container/service/slapd/assets/schema.ldif"
winpty docker exec openldap sh -c "ldapadd -Y EXTERNAL -H ldapi:/// -f /container/service/slapd/assets/01_ou_structure.ldif"
winpty docker exec openldap sh -c "ldapadd -Y EXTERNAL -H ldapi:/// -f /container/service/slapd/assets/02_groups.ldif"
winpty docker exec openldap sh -c "ldapadd -Y EXTERNAL -H ldapi:/// -f /container/service/slapd/assets/03_users.ldif"

echo " LDAP setup completed successfully."
\end{lstlisting}

\subsubsection*{Nutzen}

Durch die Verwendung dieses Skripts konnte eine stabile, wiederholbare und leicht startbare LDAP-Testumgebung geschaffen werden. Es diente während der gesamten Testphase zur Vorbereitung der Tests und ermöglichte die zuverlässige Rücksetzung des Systemzustands nach jeder Testausführung.


\subsubsection*{Beispielinhalte der LDIF-Dateien}

Die LDIF-Dateien enthalten die für die Tests notwendige LDAP-Struktur, Gruppen und Testbenutzer. Im Folgenden sind exemplarische Auszüge aus den jeweiligen Dateien dargestellt.

\paragraph{Strukturdefinition - \texttt{01\_ou\_structure.ldif}}

Diese Datei definiert die grundlegenden Organisationseinheiten innerhalb des LDAP-Verzeichnisses, z.\,B. für Benutzer oder Gruppen:

\begin{lstlisting}[language=bash, caption={Auszug aus \texttt{01\_ou\_structure.ldif}}]
dn: ou=people,dc=paxportal,dc=ch
objectClass: organizationalUnit
ou: people

dn: ou=groups,dc=paxportal,dc=ch
objectClass: organizationalUnit
ou: groups
\end{lstlisting}

\paragraph{Gruppendefinition - \texttt{02\_groups.ldif}}

Diese Datei legt LDAP-Gruppen an, denen Benutzer später zugewiesen werden:

\begin{lstlisting}[language=bash, caption={Auszug aus \texttt{02\_groups.ldif}}]
dn: cn=BROKERPORTAL,ou=groups,dc=paxportal,dc=ch
objectClass: groupOfNames
cn: BROKERPORTAL
member: cn=placeholder,dc=paxportal,dc=ch
\end{lstlisting}

\paragraph{Testbenutzer - \texttt{03\_users.ldif}}

Diese Datei enthält vorab definierte Testnutzer, die für Integrationstests verwendet werden:

\begin{lstlisting}[language=bash, caption={Auszug aus \texttt{03\_users.ldif}}]
dn: cn=Marco Frei,ou=people,dc=paxportal,dc=ch
objectClass: inetOrgPerson
objectClass: paxPeople
cn: Marco Frei
sn: Frei
givenName: Marco
mail: marco.frei@pax.ch
uid: 100
\end{lstlisting}

\subsubsection*{Zusammenfassung}

Durch die Kombination aus Container-Neustart und gezieltem Import der LDIF-Dateien wird die Testumgebung für jede Ausführung in einen sauberen Ausgangszustand zurückgesetzt. Die Strukturdateien sind modular aufgebaut und ermöglichen eine klare Trennung zwischen OUs, Gruppen und Benutzern.
